% Chương 5 (sửa trực tiếp; bản tiếng Việt)
\chapter{KẾT LUẬN VÀ HƯỚNG PHÁT TRIỂN}

\section{Trả lời các câu hỏi nghiên cứu}

\textbf{RQ1, tính khả thi trên thiết bị biên: có.} Trên Raspberry Pi 5 với Hailo-8, cấu hình triển khai (mô hình n có chia ô) chạy thời gian thực trong 15 phút:

\begin{itemize}
\tightlist
\item
  30 FPS với độ trễ 47 ms
\item
  4.1 W đo được trên bo mạch, cộng khoảng 2.5 W cho bộ tăng tốc
\item
  nhiệt độ SoC cao nhất 60 °C, không bị giảm xung
\end{itemize}

Pipeline bị giới hạn bởi bộ tăng tốc, không phải CPU. Mô hình m có chia ô chỉ đạt 7.9 FPS, đủ cho phân tích nhưng không đủ cho camera trực tiếp nếu không bỏ khung hình.

\textbf{RQ2, độ bao phủ: có thể cải thiện đáng kể mà không cần huấn luyện lại.} Chia ô vùng cuối lớp nâng tỷ lệ học sinh SCB được tìm thấy từ 55\% lên 72\%, và ở một phần ba xa nhất từ 25\% lên 66\%. Kỹ thuật này tăng hơn gấp đôi số học sinh tìm được trong một phòng tự học đông đúc. Loại bỏ trùng lặp cùng bộ khung xương và neo chỗ ngồi ở phần đầu đưa sĩ số suy ra đến gần sĩ số thực ở buổi giảng (26 cho 28 học sinh), nhưng vẫn đếm thừa ở phòng thi (29 cho 23 học viên và một giám thị) và ở phòng tự học đông đúc (59 cho khoảng 50).

Trong một phòng nhỏ quay ở khoảng cách gần (clip\_5), số chỗ ngồi cũng bị nhân lên, thành 13-16 cho 10 học sinh, nên ước lượng sĩ số chỉ là con số tham khảo, không phải số người chính xác.

\textbf{RQ3, quyết định cho từng học sinh: hiệu chỉnh tốt hơn độ chính xác trên benchmark.} Các bộ phân loại huấn luyện trên SCB xếp hạng tốt các hành động của SCB (AUC trong ảnh 0.83) nhưng chuyển giao kém:

\begin{itemize}
\tightlist
\item
  \textbf{Video buổi giảng:} chúng gọi 50-82\% của một buổi giảng chăm chú là mất tập trung.
\item
  \textbf{Nhãn người gán:} chúng đánh dấu sai 84-159 trên 262 thời điểm tập trung.
\item
  \textbf{Nguyên nhân:} SCB không có nhãn ``nghe giảng'', nên chúng học được rằng ngẩng đầu nghĩa là mất tập trung.
\item
  \textbf{Bộ dữ liệu thứ hai:} bổ sung ClaBehavior, có lớp ngẩng đầu đó, chỉ hạ tỷ lệ trên buổi giảng xuống 44-51\% với cây quyết định và MLP, và một mô hình huấn luyện trên bộ dữ liệu này chỉ đạt AUC 0.70-0.81 trên bộ kia.
\end{itemize}

Các luật đã hiệu chỉnh có rất ít báo động nhầm (độ đặc hiệu 95\% trên nhãn người gán ở cấu hình triển khai; 1-2\% trên buổi giảng). Ban đầu chúng bỏ sót phần lớn hành vi mất tập trung thực (11\% trên buổi dàn dựng), vì góc yaw của đầu dựa trên tai sụp đổ khi nhìn nghiêng. Một thước đo góc quay đầu thứ hai, tương đối so với cơ thể, nâng con số này lên 56\% với cả hai mô hình, không thêm báo động nhầm nào ở đó, với chi phí 7-10 điểm ở phòng thi. Trên một lớp 10 học sinh được gán nhãn sau khi mọi ngưỡng đã cố định (clip\_5), các luật trong cấu hình triển khai giữ độ đặc hiệu 89\%. Việc dùng điện thoại vẫn không thể thấy được qua tư thế.

\textbf{RQ4, chỉ số cấp lớp: có, với độ rộng trung thực.} Ghi nhớ chỗ ngồi, gán trọng số theo vùng và khoảng Wilson quanh giá trị hiển thị cho ra một tỷ lệ luôn nằm trong khoảng của nó. Khi không biết sĩ số thực, khoảng khá rộng (15-44 điểm trên các video), phản ánh một mẫu nhỏ như vậy chỉ nói được ít điều. Cung cấp \texttt{-\/-class-size} sẽ thu hẹp khoảng này.

\section{Hướng phát triển}

\begin{enumerate}
\def\labelenumi{\arabic{enumi}.}
\tightlist
\item
  \textbf{Kiểm tra độc lập ở quy mô lớp học.} Chạy buổi dàn dựng 9 người theo kịch bản (research/staged/script\_v2.csv). clip\_5 đã xác nhận độ đặc hiệu của các luật trên một lớp chưa từng dùng để tinh chỉnh, nhưng recall vẫn chỉ dựa trên một buổi hai người.
\item
  \textbf{Nhãn huấn luyện đặc thù cho lớp học.} ClaBehavior cho thấy một bộ dữ liệu công khai có lớp ngẩng đầu là chưa đủ. Các công cụ gán nhãn và dàn dựng có thể xây dựng một bộ dữ liệu từ chính các lớp học mục tiêu, với học sinh ngẩng đầu chú ý được gán nhãn theo thời gian. Một bộ quyết định học máy huấn luyện trên bộ dữ liệu đó, dùng giá trị gốc theo từng chỗ ngồi làm đặc trưng, có thể kết hợp độ nhạy của các bộ phân loại với độ đặc hiệu của các luật.
\item
  \textbf{Việc dùng điện thoại.} Phát hiện điện thoại cần một bộ phát hiện đối tượng. Trên Hailo-8, nó có thể chạy trên cùng các khung hình, đổi lại tốc độ khung hình giảm. Chỉ một cờ ``thấy điện thoại'' theo từng chỗ ngồi được phép rời khỏi thiết bị.
\item
  \textbf{Triển khai thực địa dài hơn.} Vài tuần sử dụng sẽ cho thấy độ trôi của hiệu chỉnh, ảnh hưởng của ánh sáng và việc đổi chỗ ngồi, và liệu giáo viên có thấy chỉ số hữu ích hay không. Việc này cần một nghiên cứu có sự đồng ý, được thiết kế cùng nhà trường.
\item
  \textbf{Chạy trực tiếp m có chia ô.} Một hàng đợi có bỏ khung hình hoặc giảm tốc độ camera sẽ giúp cấu hình chính xác nhất chạy trực tiếp mà độ trễ không tăng dần.
\item
  \textbf{Hình học dựa trên điểm khóa.} Neo chỗ ngồi theo vai và đo độ xoay thân so với phần thân trên nhìn thấy được thay vì so với hộp. Điều này loại bỏ hai lỗi thấy ở khoảng cách gần trong clip\_5: số chỗ ngồi bị nhân lên và hộp toàn thân bị hiểu là thân xoay.
\item
  \textbf{Đánh giá từ phía giáo viên.} Phép thử thực sự của một chỉ số cấp lớp là nó có giúp ích cho việc dạy hay không, điều này cần một nghiên cứu người dùng chứ không phải thêm các độ đo độ chính xác.
\end{enumerate}
